\documentclass[11pt,a4paper]{article}
\usepackage[utf8]{inputenc}
\usepackage[margin=1in]{geometry}
\usepackage{amsmath}
\usepackage{amssymb}
\usepackage{listings}
\usepackage{xcolor}
\usepackage{hyperref}
\usepackage{graphicx}
\usepackage{tikz}
\usepackage{float}

% Code listing style
\lstset{
    basicstyle=\ttfamily\small,
    keywordstyle=\color{blue},
    commentstyle=\color{green!60!black},
    stringstyle=\color{red},
    breaklines=true,
    frame=single,
    numbers=left,
    numberstyle=\tiny\color{gray},
    language=Python
}

\title{Comprehensive Notes: Cart-Pendulum with 2-DOF Planar Manipulator\\
\large Coupled Multi-Robot System in NVIDIA Isaac Sim}
\author{Script: \texttt{script\_cart\_pendulum\_manipulator\_basic\_run.py}}
\date{\today}

\begin{document}

\maketitle

\tableofcontents
\newpage

\section{System Overview}

This document provides comprehensive technical documentation for a coupled multi-robot system implemented in NVIDIA Isaac Sim. The system consists of two independent robots connected via a coupling joint:

\begin{enumerate}
    \item \textbf{Cart-Pendulum System}: A cart moving on rails with a hanging pendulum
    \item \textbf{2-DOF Planar Manipulator}: A two-link robotic arm operating in the XY plane
    \item \textbf{Coupling Joint}: Fixed joint connecting the manipulator end-effector to the cart edge
\end{enumerate}

\subsection{Educational Objectives}

\begin{itemize}
    \item Understand forward and inverse kinematics for planar manipulators
    \item Learn Jacobian computation and differential inverse kinematics
    \item Study multi-body dynamics with mechanical constraints
    \item Analyze force transmission through coupling joints
    \item Implement trajectory planning and control
    \item Explore object-oriented design patterns for robotics
\end{itemize}

\subsection{System Dynamics}

The coupled system exhibits complex dynamics where:
\begin{itemize}
    \item Manipulator motion drives cart translation along the X-axis
    \item Cart acceleration induces pendulum swing (inertial forces)
    \item Coupling joint transmits forces/torques between subsystems
    \item System demonstrates multi-body constrained dynamics
\end{itemize}

\section{Application Initialization}

\subsection{Isaac Sim Launch}

\begin{lstlisting}
simulation_app = SimulationApp({
    "headless": False,
    "width": 1280,
    "height": 720,
})
\end{lstlisting}

\textbf{Critical Note:} SimulationApp \textit{must} be created before importing any Isaac Sim modules. This initializes the Omniverse Kit framework.

\textbf{Configuration Parameters:}
\begin{itemize}
    \item \texttt{headless=False}: Enable GUI window (set True for batch/server runs)
    \item \texttt{width=1280, height=720}: Window resolution (affects rendering quality)
\end{itemize}

\subsection{Module Imports}

\begin{lstlisting}
# USD (Universal Scene Description)
from pxr import UsdGeom, UsdLux, Gf, UsdShade, Sdf

# Isaac Sim Core
from omni.isaac.core import World
from isaacsim.core.experimental.prims import Articulation, RigidPrim

# URDF Import
from isaacsim.asset.importer.urdf import _urdf

# Math Libraries
import numpy as np
import torch
import warp as wp
\end{lstlisting}

\textbf{Module Categories:}
\begin{itemize}
    \item \textbf{USD/Pixar}: Scene graph manipulation (geometry, transforms, materials)
    \item \textbf{Physics}: Articulation (multi-link robots), rigid body dynamics
    \item \textbf{URDF}: Robot description format conversion
    \item \textbf{Math}: NumPy (CPU arrays), PyTorch (GPU tensors), Warp (parallel computing)
\end{itemize}

\section{URDF to USD Conversion}

\subsection{Why Convert URDF to USD?}

\begin{itemize}
    \item \textbf{URDF} (Unified Robot Description Format): XML-based, ROS standard, human-readable
    \item \textbf{USD} (Universal Scene Description): Binary format, Pixar standard, optimized for real-time rendering and physics
    \item Isaac Sim's physics engine requires USD format for simulation
\end{itemize}

\subsection{Conversion Function}

\begin{lstlisting}
def convert_urdf_to_usd(urdf_path, output_usd_path, import_config=None):
    urdf_interface = _urdf.acquire_urdf_interface()
    
    config = _urdf.ImportConfig()
    config.convex_decomp = import_config.get("convex_decomp", False)
    config.fix_base = import_config.get("fix_base", False)
    config.make_default_prim = import_config.get("make_default_prim", True)
    config.self_collision = import_config.get("self_collision", False)
    config.distance_scale = import_config.get("distance_scale", 1.0)
    config.density = import_config.get("density", 0.0)
    
    result, prim_path = omni.kit.commands.execute(
        "URDFParseAndImportFile",
        urdf_path=urdf_path,
        import_config=config,
        dest_path=output_usd_path
    )
    
    return prim_path
\end{lstlisting}

\textbf{Import Configuration:}
\begin{itemize}
    \item \texttt{convex\_decomp}: Decompose complex meshes into convex hulls for collision
    \item \texttt{fix\_base}: Fix base link to world (no floating base)
    \item \texttt{make\_default\_prim}: Set as default USD prim for stage
    \item \texttt{self\_collision}: Enable collision detection between robot links
    \item \texttt{distance\_scale}: Unit conversion factor (default 1.0 = meters)
    \item \texttt{density}: Override material density (0.0 = use URDF values)
\end{itemize}

\subsection{Automatic Regeneration}

\begin{lstlisting}
def needs_regeneration(urdf_path, usd_path):
    if not os.path.exists(usd_path):
        return True  # USD doesn't exist
    
    urdf_mtime = os.path.getmtime(urdf_path)
    usd_mtime = os.path.getmtime(usd_path)
    
    return urdf_mtime > usd_mtime  # Regenerate if URDF is newer
\end{lstlisting}

\textbf{Purpose:} Only regenerate USD if URDF file has been modified, saving time during development.

\section{Object-Oriented Architecture}

\subsection{Design Patterns}

The implementation uses several software design patterns:

\begin{enumerate}
    \item \textbf{Template Method Pattern} (RobotBase): Define skeleton, subclasses implement specifics
    \item \textbf{Facade Pattern} (SceneManager): Simple interface to complex subsystems
    \item \textbf{Dataclass Pattern} (RobotParams): Clean configuration containers
\end{enumerate}

\subsection{Class Hierarchy}

\begin{verbatim}
RobotBase (Abstract Base Class)
|-- CartPendulum
`-- PlanarManipulator

SceneManager (Orchestrator)
|-- CartPendulum instance
|-- PlanarManipulator instance
|-- World (physics)
`-- USD Stage (scene)
\end{verbatim}

\subsection{RobotBase Abstract Class}

\begin{lstlisting}
class RobotBase(ABC):
    def __init__(self, params: RobotParams):
        self.params = params
        self.state: Optional[RobotState] = None
    
    @abstractmethod
    def set_joint_properties(self):
        """Must be implemented by subclasses"""
        pass
    
    def load_to_stage(self, stage):
        """Shared implementation for all robots"""
        prim = stage.OverridePrim(self.params.prim_path)
        prim.GetReferences().AddReference(self.params.usd_path)
    
    def initialize_articulation(self):
        """Create Articulation after world.reset()"""
        robot = Articulation(self.params.prim_path)
        self.state = RobotState(
            robot=robot,
            num_dof=robot.num_dofs,
            dof_names=robot.dof_names
        )
\end{lstlisting}

\textbf{Key Methods:}
\begin{itemize}
    \item \texttt{load\_to\_stage()}: Add robot USD to scene graph
    \item \texttt{initialize\_articulation()}: Create physics object after reset
    \item \texttt{set\_joint\_properties()}: Configure damping, friction (abstract)
    \item \texttt{set\_joint\_positions()}: Position control
    \item \texttt{get\_joint\_positions()}: State query
\end{itemize}

\section{Forward Kinematics}

\subsection{Problem Statement}

\textbf{Given:} Joint angles $\theta_1, \theta_2$

\textbf{Find:} End-effector position $(x, y, z)$ in world coordinates

\subsection{Mathematical Derivation}

For a 2-DOF planar manipulator with link lengths $L_1$ and $L_2$:

\subsubsection{Position in Base Frame}

The end-effector position relative to the base is:
\begin{align}
x_{\text{base}} &= L_1 \cos(\theta_1) + L_2 \cos(\theta_1 + \theta_2) \\
y_{\text{base}} &= L_1 \sin(\theta_1) + L_2 \sin(\theta_1 + \theta_2)
\end{align}

\textbf{Intuition:}
\begin{itemize}
    \item First link contributes: $(L_1 \cos\theta_1, L_1 \sin\theta_1)$
    \item Second link rotates by $(\theta_1 + \theta_2)$ absolute angle
    \item Second link contributes: $(L_2 \cos(\theta_1+\theta_2), L_2 \sin(\theta_1+\theta_2))$
    \item Sum gives total position
\end{itemize}

\subsubsection{Transform to World Frame}

\begin{align}
x_{\text{world}} &= x_{\text{base}} + x_{\text{base\_offset}} \\
y_{\text{world}} &= y_{\text{base}} + y_{\text{base\_offset}} \\
z_{\text{world}} &= z_{\text{base}} + z_{\text{base\_offset}}
\end{align}

\subsection{Implementation}

\begin{lstlisting}
def forward_kinematics(self, theta1, theta2):
    L1 = self.link_lengths[0]
    L2 = self.link_lengths[1]
    
    # Base frame position
    x_rel = L1 * np.cos(theta1) + L2 * np.cos(theta1 + theta2)
    y_rel = L1 * np.sin(theta1) + L2 * np.sin(theta1 + theta2)
    
    # Transform to world coordinates
    x_world = self.params.position[0] + x_rel
    y_world = self.params.position[1] + y_rel
    z_world = self.params.position[2] + 1.3  # Base height
    
    return (x_world, y_world, z_world)
\end{lstlisting}

\subsection{Properties}

\begin{itemize}
    \item \textbf{Uniqueness}: Forward kinematics has a unique solution
    \item \textbf{Computational Cost}: $O(1)$ - just trigonometric functions
    \item \textbf{Workspace}: Annulus with inner radius $|L_1-L_2|$, outer radius $L_1+L_2$
\end{itemize}

\section{Inverse Kinematics}

\subsection{Problem Statement}

\textbf{Given:} Desired end-effector position $(x_{\text{target}}, y_{\text{target}}, z_{\text{target}})$

\textbf{Find:} Joint angles $(\theta_1, \theta_2)$ that achieve this position

\subsection{Analytical Solution}

\subsubsection{Step 1: Transform to Base Frame}

Convert world coordinates to base-relative coordinates using transformation matrix:
\begin{equation}
\begin{bmatrix} x_{\text{rel}} \\ y_{\text{rel}} \\ z_{\text{rel}} \end{bmatrix} = T_{\text{world}\to\text{base}} \begin{bmatrix} x_{\text{target}} \\ y_{\text{target}} \\ z_{\text{target}} \end{bmatrix}
\end{equation}

\subsubsection{Step 2: Check Reachability}

Compute distance from base to target:
\begin{equation}
r = \sqrt{x_{\text{rel}}^2 + y_{\text{rel}}^2}
\end{equation}

\textbf{Reachability condition:}
\begin{equation}
|L_1 - L_2| \leq r \leq L_1 + L_2
\end{equation}

If violated, target is unreachable (too far or too close).

\subsubsection{Step 3: Compute Joint 2 Angle}

Using the Law of Cosines:
\begin{equation}
\cos(\theta_2) = \frac{r^2 - L_1^2 - L_2^2}{2 L_1 L_2}
\end{equation}

For \textbf{elbow-down} configuration (negative $\theta_2$):
\begin{equation}
\theta_2 = -\arccos\left(\frac{r^2 - L_1^2 - L_2^2}{2 L_1 L_2}\right)
\end{equation}

\subsubsection{Step 4: Compute Joint 1 Angle}

Angle to target from base:
\begin{equation}
\phi = \arctan2(y_{\text{rel}}, x_{\text{rel}})
\end{equation}

Geometric offset angle:
\begin{equation}
\alpha = \arctan2(L_2 \sin\theta_2, L_1 + L_2 \cos\theta_2)
\end{equation}

Final joint 1 angle:
\begin{equation}
\theta_1 = \phi - \alpha
\end{equation}

\subsection{Implementation}

\begin{lstlisting}
def inverse_kinematics(self, target_x, target_y, target_z):
    # Transform to base frame
    base_coords = self.transform_point_world_to_base(
        target_x, target_y, target_z)
    x_rel, y_rel, z_rel = base_coords
    
    # Check reachability
    r = np.sqrt(x_rel**2 + y_rel**2)
    if r > (L1 + L2) or r < abs(L1 - L2):
        return None  # Unreachable
    
    # Law of cosines for theta2
    cos_theta2 = (r**2 - L1**2 - L2**2) / (2 * L1 * L2)
    cos_theta2 = np.clip(cos_theta2, -1.0, 1.0)
    theta2 = -np.arccos(cos_theta2)  # Elbow-down
    
    # Solve for theta1
    phi = np.arctan2(y_rel, x_rel)
    alpha = np.arctan2(L2 * np.sin(theta2), L1 + L2 * np.cos(theta2))
    theta1 = phi - alpha
    
    return (theta1, theta2)
\end{lstlisting}

\subsection{Multiple Solutions}

For a 2-DOF planar manipulator, there are typically \textbf{two solutions}:
\begin{itemize}
    \item \textbf{Elbow-down}: $\theta_2 < 0$ (implemented above)
    \item \textbf{Elbow-up}: $\theta_2 > 0$ (use $+\arccos$ instead)
\end{itemize}

The choice depends on application constraints (obstacle avoidance, joint limits, etc.).

\section{Jacobian Matrix}

\subsection{Definition}

The \textbf{Jacobian matrix} $J$ relates joint velocities to end-effector velocities:
\begin{equation}
\begin{bmatrix} v_x \\ v_y \end{bmatrix} = J \begin{bmatrix} \dot{\theta}_1 \\ \dot{\theta}_2 \end{bmatrix}
\end{equation}

where $v_x, v_y$ are end-effector linear velocities and $\dot{\theta}_1, \dot{\theta}_2$ are joint angular velocities.

\subsection{Mathematical Derivation}

Starting from forward kinematics:
\begin{align}
x &= L_1 \cos\theta_1 + L_2 \cos(\theta_1 + \theta_2) \\
y &= L_1 \sin\theta_1 + L_2 \sin(\theta_1 + \theta_2)
\end{align}

Taking partial derivatives:
\begin{align}
\frac{\partial x}{\partial \theta_1} &= -L_1 \sin\theta_1 - L_2 \sin(\theta_1 + \theta_2) \\
\frac{\partial x}{\partial \theta_2} &= -L_2 \sin(\theta_1 + \theta_2) \\
\frac{\partial y}{\partial \theta_1} &= L_1 \cos\theta_1 + L_2 \cos(\theta_1 + \theta_2) \\
\frac{\partial y}{\partial \theta_2} &= L_2 \cos(\theta_1 + \theta_2)
\end{align}

\subsection{Jacobian Matrix}

\begin{equation}
J(\theta_1, \theta_2) = \begin{bmatrix}
-L_1 \sin\theta_1 - L_2 \sin(\theta_1+\theta_2) & -L_2 \sin(\theta_1+\theta_2) \\
L_1 \cos\theta_1 + L_2 \cos(\theta_1+\theta_2) & L_2 \cos(\theta_1+\theta_2)
\end{bmatrix}
\end{equation}

\subsection{Implementation}

\begin{lstlisting}
def compute_jacobian(self, theta1, theta2):
    L1 = self.link_lengths[0]
    L2 = self.link_lengths[1]
    
    s1 = np.sin(theta1)
    c1 = np.cos(theta1)
    s12 = np.sin(theta1 + theta2)
    c12 = np.cos(theta1 + theta2)
    
    J = np.array([
        [-L1*s1 - L2*s12, -L2*s12],
        [ L1*c1 + L2*c12,  L2*c12]
    ])
    
    return J
\end{lstlisting}

\subsection{Singularities}

\textbf{Singularity:} Configuration where $\det(J) = 0$, meaning robot loses degrees of freedom.

For 2-DOF planar manipulator, singularities occur when:
\begin{equation}
\det(J) = L_1 L_2 \sin\theta_2 = 0
\end{equation}

This happens when $\theta_2 = 0$ or $\theta_2 = \pi$ (arm fully extended or folded).

\subsection{Condition Number}

\textbf{Condition number:} Measure of numerical stability
\begin{equation}
\kappa(J) = \|J\| \cdot \|J^{-1}\|
\end{equation}

\begin{itemize}
    \item $\kappa < 10$: Well-conditioned (good for control)
    \item $10 < \kappa < 100$: Moderately conditioned
    \item $\kappa > 100$: Poorly conditioned (near singularity)
\end{itemize}

\section{Differential Inverse Kinematics}

\subsection{Concept}

Instead of computing exact joint angles (analytical IK), \textbf{differential IK} computes incremental changes $\Delta\theta$ to move end-effector toward goal.

\textbf{Algorithm:}
\begin{enumerate}
    \item Compute position error: $e = x_{\text{goal}} - x_{\text{current}}$
    \item Compute Jacobian: $J = J(\theta_{\text{current}})$
    \item Solve for joint velocity: $\dot{\theta} = J^{-1} e$
    \item Update: $\theta_{\text{new}} = \theta_{\text{old}} + \alpha \dot{\theta}$ ($\alpha$ = step size)
    \item Repeat until $\|e\| < \epsilon$
\end{enumerate}

\subsection{Methods for Computing $J^{-1}$}

\subsubsection{1. Pseudoinverse (Moore-Penrose)}

\begin{equation}
J^+ = (J^T J)^{-1} J^T
\end{equation}

\textbf{Properties:}
\begin{itemize}
    \item Minimizes $\|\Delta\theta\|$ (smallest joint motion)
    \item Exact solution for non-singular $J$
    \item \textbf{Problem:} Unstable near singularities
\end{itemize}

\subsubsection{2. Transpose Method}

\begin{equation}
\Delta\theta = k \cdot J^T e
\end{equation}

\textbf{Properties:}
\begin{itemize}
    \item Computationally cheap (no matrix inversion)
    \item Stable near singularities
    \item \textbf{Problem:} Does not minimize error exactly, slower convergence
\end{itemize}

\subsubsection{3. Damped Least Squares (Levenberg-Marquardt)}

\begin{equation}
\Delta\theta = J^T (J J^T + \lambda^2 I)^{-1} e
\end{equation}

\textbf{Properties:}
\begin{itemize}
    \item Adds damping term $\lambda$ to avoid singularities
    \item Trade-off: accuracy vs. stability (controlled by $\lambda$)
    \item \textbf{Recommended for most applications}
    \item Typical values: $\lambda \in [0.01, 0.1]$
\end{itemize}

\subsubsection{4. Singular Value Decomposition (SVD)}

\begin{align}
J &= U \Sigma V^T \\
J^+ &= V \Sigma^+ U^T
\end{align}

where $\Sigma^+$ inverts non-zero singular values and zeros out small ones.

\textbf{Properties:}
\begin{itemize}
    \item Most robust method
    \item Can filter small singular values (threshold)
    \item \textbf{Problem:} Computationally expensive
\end{itemize}

\subsection{PyTorch Implementation}

\begin{lstlisting}
def inverse_kinematics_differential(
    self, jacobian_end_effector, current_position, 
    goal_position, method="damped-least-squares", method_cfg=None
):
    if method_cfg is None:
        method_cfg = {"scale": 1.0, "damping": 0.05}
    
    scale = method_cfg.get("scale", 1.0)
    damping = method_cfg.get("damping", 0.05)
    
    # Compute error
    error = (goal_position - current_position).unsqueeze(-1)
    
    if method == "damped-least-squares":
        J_T = torch.transpose(jacobian_end_effector, 1, 2)
        lambda_I = torch.eye(2) * (damping ** 2)
        delta_theta = (scale * J_T @ torch.inverse(
            jacobian_end_effector @ J_T + lambda_I) @ error)
    
    elif method == "pseudoinverse":
        J_pinv = torch.linalg.pinv(jacobian_end_effector)
        delta_theta = (scale * J_pinv @ error)
    
    elif method == "transpose":
        J_T = torch.transpose(jacobian_end_effector, 1, 2)
        delta_theta = (scale * J_T @ error)
    
    elif method == "singular-value-decomposition":
        U, S, Vh = torch.linalg.svd(jacobian_end_effector)
        min_sv = method_cfg.get("min_singular_value", 1e-5)
        inv_s = torch.where(S > min_sv, 1.0 / S, torch.zeros_like(S))
        J_pinv = Vh.transpose(1,2) @ torch.diag_embed(inv_s) @ U.transpose(1,2)
        delta_theta = (scale * J_pinv @ error)
    
    return delta_theta.squeeze(-1)
\end{lstlisting}

\section{Joint Properties}

\subsection{Damping}

\textbf{Physical model:} Viscous damping opposes motion proportional to velocity
\begin{equation}
\tau_{\text{damp}} = -b \dot{\theta}
\end{equation}

where $b$ is the damping coefficient (units: N·m·s/rad).

\textbf{Effects:}
\begin{itemize}
    \item Dissipates energy from system
    \item Prevents oscillations
    \item Stabilizes motion
\end{itemize}

\textbf{Typical values:}
\begin{itemize}
    \item Cart slider: $b = 0.5$ N·m·s/rad (moderate damping)
    \item Pendulum: $b = 0.01$ N·m·s/rad (low damping for natural swing)
    \item Manipulator joints: $b = 0.1$ N·m·s/rad (light damping)
\end{itemize}

\subsection{Stiffness}

\textbf{Physical model:} Spring restoring force proportional to displacement
\begin{equation}
\tau_{\text{spring}} = -k (\theta - \theta_{\text{rest}})
\end{equation}

where $k$ is stiffness coefficient (units: N·m/rad).

\textbf{Effects:}
\begin{itemize}
    \item Returns joint to rest position
    \item Creates passive compliance
    \item Can stabilize equilibrium
\end{itemize}

\textbf{Configuration in this system:}
\begin{itemize}
    \item Cart slider: $k = 625$ N·m/rad (prevents drift)
    \item Pendulum: $k = 0.1$ N·m/rad (near-zero for free motion)
    \item Manipulator: $k = 0$ N·m/rad (no restoring force)
\end{itemize}

\subsection{Friction}

\textbf{Physical model:} Opposes motion independent of velocity
\begin{equation}
\tau_{\text{friction}} = -\mu \cdot \text{sign}(\dot{\theta})
\end{equation}

where $\mu$ is friction coefficient.

\textbf{Types:}
\begin{itemize}
    \item \textbf{Static friction}: Prevents motion when at rest
    \item \textbf{Dynamic friction}: Opposes motion when moving
\end{itemize}

\textbf{Configuration:}
\begin{itemize}
    \item Cart slider: $\mu = 0.1$ (low friction for smooth motion)
    \item Pendulum: $\mu = 0$ (frictionless pivot)
    \item Manipulator: $\mu = 0$ (actuated joints)
\end{itemize}

\subsection{Implementation}

\begin{lstlisting}
def _set_joint_damping(self, joint_prim, damping_value, 
                      stiffness_value=0.0, joint_type="revolute"):
    from pxr import UsdPhysics
    
    if joint_type == "revolute":
        drive = UsdPhysics.DriveAPI.Apply(joint_prim, "angular")
    elif joint_type == "prismatic":
        drive = UsdPhysics.DriveAPI.Apply(joint_prim, "linear")
    
    drive.CreateTypeAttr().Set("force")
    drive.CreateDampingAttr().Set(float(damping_value))
    drive.CreateStiffnessAttr().Set(float(stiffness_value))

def _set_joint_friction(self, joint_prim, friction_value):
    from pxr import PhysxSchema
    
    physx_joint = PhysxSchema.PhysxJointAPI.Apply(joint_prim)
    friction_attr = physx_joint.CreateJointFrictionAttr()
    friction_attr.Set(float(friction_value))
\end{lstlisting}

\section{Coupling Joint}

\subsection{Purpose}

Connect manipulator end-effector to cart edge, enabling:
\begin{itemize}
    \item Force/torque transmission between subsystems
    \item Manipulator drives cart motion
    \item Cart acceleration swings pendulum
    \item Demonstrates multi-body constrained dynamics
\end{itemize}

\subsection{Joint Type: Fixed Joint}

\textbf{Fixed joint:} Rigid connection with 6 constraints (no relative motion)
\begin{itemize}
    \item Constrains: 3 translations + 3 rotations
    \item End-effector and cart edge move as single rigid body
    \item Forces transmitted directly through connection
\end{itemize}

\subsection{Implementation}

\begin{lstlisting}
def create_ee_cart_joint(self):
    from pxr import UsdPhysics
    
    ee_path = f"{self.manipulator.params.prim_path}/manipulator_link_2/manipulator_ee"
    cart_path = f"{self.cart_pendulum.params.prim_path}/cart"
    joint_path = "/World/ee_cart_coupling_joint"
    
    # Get transformations
    ee_world_transform = self.manipulator.get_world_transform(ee_prim)
    cart_world_transform = self.cart_pendulum.get_world_transform(cart_prim)
    
    # Compute joint position in cart's local frame
    cart_to_world_inv = cart_world_transform.GetInverse()
    joint_pos_in_cart = cart_to_world_inv.Transform(ee_translation)
    
    # Create fixed joint
    joint = UsdPhysics.FixedJoint.Define(stage, joint_path)
    joint.CreateBody0Rel().SetTargets([ee_path])
    joint.CreateBody1Rel().SetTargets([cart_path])
    joint.CreateLocalPos0Attr().Set(Gf.Vec3f(0, 0, 0))
    joint.CreateLocalPos1Attr().Set(Gf.Vec3f(joint_pos_in_cart))
    
    return True
\end{lstlisting}

\subsection{Force Analysis}

When manipulator moves with acceleration $\ddot{\theta}_1, \ddot{\theta}_2$:

\begin{enumerate}
    \item End-effector experiences acceleration: $\mathbf{a}_{\text{EE}} = J \ddot{\boldsymbol{\theta}} + \dot{J} \dot{\boldsymbol{\theta}}$
    \item Coupling joint transmits force to cart: $\mathbf{F}_{\text{coupling}}$
    \item Cart accelerates along rail: $\ddot{x}_{\text{cart}} = F_x / m_{\text{cart}}$
    \item Pendulum experiences inertial force: $F_{\text{inertial}} = m_{\text{pendulum}} \ddot{x}_{\text{cart}}$
    \item Pendulum swings: $\ddot{\phi}_{\text{pendulum}} = -\frac{g}{L} \sin\phi - \frac{\ddot{x}_{\text{cart}}}{L} \cos\phi$
\end{enumerate}

\section{Simulation Modes}

\subsection{Scene Visualization (scene-viz)}

\textbf{Purpose:} Static scene viewing without physics simulation

\textbf{Features:}
\begin{itemize}
    \item Display robots at initial configuration
    \item No physics computation
    \item Useful for debugging USD scene structure
\end{itemize}

\subsection{End-Effector Trajectory (ee-trajectory)}

\textbf{Purpose:} Move manipulator EE along X-axis using inverse kinematics

\textbf{Algorithm:}
\begin{enumerate}
    \item Compute initial EE position using forward kinematics
    \item Generate sinusoidal trajectory along X-axis
    \item For each waypoint, solve inverse kinematics
    \item Set joint positions and step simulation
\end{enumerate}

\textbf{Trajectory equation:}
\begin{equation}
x(t) = x_0 + A \sin\left(\frac{2\pi}{T} t\right)
\end{equation}

where $A$ is amplitude (meters), $T$ is period (seconds).

\subsection{Cart Toward Manipulator (cart-toward-manipulator)}

\textbf{Purpose:} Autonomously move cart edge to align with manipulator EE

\textbf{Control Law:}
\begin{equation}
v_{\text{cart}} = K_p (x_{\text{EE}} - x_{\text{cart\_edge}})
\end{equation}

\textbf{Convergence:} Stops when $|x_{\text{EE}} - x_{\text{cart\_edge}}| < \epsilon$ (threshold = 0.05mm)

\subsection{Coupled Motion (coupled-motion)}

\textbf{Purpose:} Demonstrate force transmission through coupling joint

\textbf{Features:}
\begin{enumerate}
    \item Initialize robots in aligned position
    \item Create fixed joint between EE and cart
    \item Generate manipulator trajectory
    \item Execute trajectory: manipulator pulls/pushes cart
    \item Record video of coupled motion
    \item Log positions: EE, cart, pendulum
\end{enumerate}

\textbf{Observations:}
\begin{itemize}
    \item Manipulator motion directly drives cart translation
    \item Cart acceleration induces pendulum swing
    \item System exhibits complex coupled dynamics
    \item Energy transfers between subsystems
\end{itemize}

\section{Video Recording}

\subsection{Setup}

\begin{lstlisting}
def start_video_recording(self, output_path, resolution=(1920, 1080), fps=60):
    viewport_api = get_active_viewport()
    base_path = output_path.replace('.mp4', '')
    
    result = omni.kit.commands.execute(
        "StartMovieCapture",
        file_path=base_path,
        resolution=resolution,
        fps=fps,
        viewport_api=viewport_api
    )
    
    self.is_recording = True
    return result
\end{lstlisting}

\textbf{Parameters:}
\begin{itemize}
    \item \texttt{output\_path}: Full path to MP4 file
    \item \texttt{resolution}: $(width, height)$ in pixels (e.g., 1920×1080 for Full HD)
    \item \texttt{fps}: Frames per second (60 for smooth motion)
\end{itemize}

\subsection{Recording Workflow}

\begin{enumerate}
    \item Start recording before simulation loop
    \item Execute simulation (render=True in each step)
    \item Stop recording after trajectory completion
    \item Video automatically saved to specified path
\end{enumerate}

\section{Educational Exercises}

\subsection{Exercise 1: Workspace Analysis}

\textbf{Task:} Compute and visualize the manipulator's reachable workspace.

\textbf{Steps:}
\begin{enumerate}
    \item Generate grid of $(\theta_1, \theta_2)$ combinations
    \item For each configuration, compute EE position using forward kinematics
    \item Plot points to visualize workspace (annulus shape)
    \item Identify inner radius: $r_{\min} = |L_1 - L_2|$
    \item Identify outer radius: $r_{\max} = L_1 + L_2$
\end{enumerate}

\subsection{Exercise 2: Singularity Analysis}

\textbf{Task:} Find and characterize singularity configurations.

\textbf{Steps:}
\begin{enumerate}
    \item Compute Jacobian for range of joint angles
    \item Calculate determinant: $\det(J) = L_1 L_2 \sin\theta_2$
    \item Identify singularities: $\theta_2 = 0$ (extended), $\theta_2 = \pi$ (folded)
    \item Compute condition number near singularities
    \item Test differential IK behavior near singularities
\end{enumerate}

\subsection{Exercise 3: Trajectory Optimization}

\textbf{Task:} Generate smooth trajectory minimizing joint velocities/accelerations.

\textbf{Approach:}
\begin{enumerate}
    \item Define waypoints in task space
    \item Solve inverse kinematics for each waypoint
    \item Fit smooth spline through joint-space waypoints
    \item Compute velocity/acceleration profiles
    \item Evaluate smoothness metrics (jerk, etc.)
\end{enumerate}

\subsection{Exercise 4: Pendulum Swing Analysis}

\textbf{Task:} Analyze pendulum response to cart motion.

\textbf{Analysis:}
\begin{enumerate}
    \item Record cart position $x_{\text{cart}}(t)$ during coupled motion
    \item Record pendulum angle $\phi(t)$
    \item Compute cart acceleration: $\ddot{x}_{\text{cart}}(t)$
    \item Verify pendulum equation: $\ddot{\phi} = -\frac{g}{L}\sin\phi - \frac{\ddot{x}_{\text{cart}}}{L}\cos\phi$
    \item Identify resonant frequencies
\end{enumerate}

\section{Troubleshooting}

\subsection{Common Issues}

\subsubsection{Inverse Kinematics Fails}

\textbf{Symptoms:} Returns None for valid-looking targets

\textbf{Causes:}
\begin{itemize}
    \item Target outside workspace (check reachability condition)
    \item Numerical precision issues (use \texttt{np.clip} for arc-cosine)
    \item Wrong coordinate frame (world vs. base)
\end{itemize}

\textbf{Solutions:}
\begin{itemize}
    \item Verify workspace boundaries: $|L_1-L_2| \leq r \leq L_1+L_2$
    \item Add tolerance to reachability check
    \item Print intermediate values for debugging
\end{itemize}

\subsubsection{Robots Not Moving}

\textbf{Symptoms:} Joint commands issued but robots remain static

\textbf{Causes:}
\begin{itemize}
    \item Articulation not initialized (call after \texttt{world.reset()})
    \item High joint stiffness/damping
    \item Physics not stepping (missing \texttt{world.step()})
\end{itemize}

\textbf{Solutions:}
\begin{itemize}
    \item Ensure initialization order: reset → initialize → set properties
    \item Check joint property values
    \item Verify simulation loop calls \texttt{world.step(render=True)}
\end{itemize}

\subsubsection{Unstable Coupling Joint}

\textbf{Symptoms:} Excessive forces, explosion, divergence

\textbf{Causes:}
\begin{itemize}
    \item Incorrect joint positioning
    \item Conflicting constraints
    \item Time step too large
\end{itemize}

\textbf{Solutions:}
\begin{itemize}
    \item Verify local positions in joint definition
    \item Use smaller time step
    \item Add compliance (use revolute instead of fixed joint)
\end{itemize}

\section{Extensions and Future Work}

\subsection{3-DOF Spatial Manipulator}

Extend to 3D workspace:
\begin{itemize}
    \item Add third revolute joint for vertical motion
    \item Implement 3D forward/inverse kinematics
    \item Compute 3×3 Jacobian matrix
\end{itemize}

\subsection{Redundancy Resolution}

For manipulators with $n > m$ DOF (redundant):
\begin{itemize}
    \item Primary task: Reach target position
    \item Secondary task: Optimize posture (joint limits, obstacle avoidance)
    \item Use null-space projection: $\Delta\theta = J^+ e + (I - J^+ J) \theta_0$
\end{itemize}

\subsection{Compliant Coupling}

Replace fixed joint with compliant spring-damper:
\begin{itemize}
    \item Add prismatic joint with spring stiffness
    \item Study vibration transmission
    \item Analyze resonances and damping
\end{itemize}

\subsection{Model Predictive Control}

Implement MPC for trajectory tracking:
\begin{itemize}
    \item Predict future states using dynamics model
    \item Optimize control inputs over horizon
    \item Account for constraints (joint limits, collision)
\end{itemize}

\section{Conclusion}

This system demonstrates fundamental robotics concepts:
\begin{itemize}
    \item Kinematics (forward, inverse, differential)
    \item Dynamics (multi-body, constraints, coupling)
    \item Control (position, trajectory tracking)
    \item Software design (OOP, design patterns)
\end{itemize}

The implementation provides a foundation for studying more complex systems and control strategies in robotics.

\section{References}

\begin{enumerate}
    \item J.J. Craig, \textit{Introduction to Robotics: Mechanics and Control}, 4th ed., Pearson, 2017
    \item B. Siciliano, O. Khatib, \textit{Springer Handbook of Robotics}, 2nd ed., Springer, 2016
    \item R.M. Murray, Z. Li, S.S. Sastry, \textit{A Mathematical Introduction to Robotic Manipulation}, CRC Press, 1994
    \item NVIDIA Isaac Sim Documentation: \url{https://docs.omniverse.nvidia.com/isaacsim/}
    \item USD Documentation: \url{https://graphics.pixar.com/usd/docs/index.html}
\end{enumerate}

\end{document}
